% Generated by tools/textbook/export_listings.js from tools/decks/mastery.js and tools/concepts/registry.json: do not edit.
\begin{chapstatus}{Stato di apprendimento di questo capitolo}
\masterychip{now}{DA COMPRENDERE OGGI}\enspace il metodo: prima il programma; le regole del corso.\par\smallskip
\masterychip{use}{DA UTILIZZARE OGGI}\enspace il programma completo, gli strumenti di PRG32 (si usa, non si studia, in questo corso).\par\smallskip
\masterychip{later}{DA COMPRENDERE PIÙ AVANTI}\enspace i componenti di PRG32 (Capitolo~\ref{ch:setup}); RISC e un ISA aperto (Capitolo~\ref{ch:history}); il progetto Verm1 (Capitolo~\ref{ch:verm1p1}).
\end{chapstatus}
\begin{chapstatus}{Percorso del capitolo}
Alla fine del capitolo saprete:
\begin{enumerate}[nosep,leftmargin=2em]
\item descrivere i quattro passi della spirale di apprendimento e dire perché la previsione viene prima dell'esecuzione
\item enunciare le regole di valutazione: una prova intercorso, due parti d'esame da superare entrambe, nessun dispositivo e nessuna IA
\end{enumerate}
Si appoggia su: niente: è il primo capitolo. Prossimo passo: Capitolo~\ref{ch:setup}.
\end{chapstatus}
